% Introduction fragment for an AASTeX manuscript.
% Citation keys match sky_subtraction.bib.
% In the parent manuscript use:
% \bibliographystyle{aasjournal}
% \bibliography{sky_subtraction}
% Add the present paper's specific motivation and contribution after this text.

\section{Introduction}
\label{sec:introduction}

Optical spectroscopy of faint sources depends on accurate sky subtraction,
because small background errors can hide spectral features or bias the
science signal. Subtracting the sky takes both an estimate of the incident
background and a description of how the instrument records it
\citep{Wyse1992}. The goal is to remove the mean sky signal without biasing
the source spectrum; the photon noise the sky contributes remains.

In long-slit spectroscopy, the simplest method interpolates the background
beneath the target from source-free regions of the slit. This samples the sky
simultaneously, but it fails when an extended source fills the slit.
Separate (offset) sky exposures avoid contamination by the source, at the
price of sensitivity to temporal variability. Nod-and-shuffle observations
alternate object and sky on the same detector pixels, which reduces response
mismatches at a cost in observing efficiency \citep{Glazebrook2001}.
Spectral sampling matters as well: \citet{Kelson2003} used the subpixel
sampling that spectral curvature provides to fit B-spline sky models on the
original detector grid, before rectification introduces interpolation
errors.
For extended, low-surface-brightness targets, variations of the instrumental line-spread function (LSF) along the slit or across the IFU field of view can become particularly important.
\citet{Katkov2011} developed a Fourier-space reconstruction technique that explicitly accounts for such LSF variations, improving sky subtraction in long-slit observations of galaxy outskirts.
Modern slit-reduction software such as \texttt{PypeIt} builds
sophisticated background modelling into the reduction workflow
\citep{Prochaska2020}.

Fiber spectroscopy instead samples the sky through dedicated fibers placed
away from the targets. Transferring that estimate to each science fiber
requires calibrating the relative throughput, the wavelength scale and,
crucially, the line-spread function (LSF). \citet{Lissandrini1994} showed
the value of night-sky emission lines for throughput calibration, together
with corrections for diffuse light and image distortions. Even so, a small
wavelength offset or an LSF mismatch can leave large residuals in and around
the sky lines. A more general approach models the two-dimensional response of
the spectrograph and detector during extraction, which retains the resolution
information needed to compare a common sky spectrum with each fiber
\citep{Bolton2010}. The DESI pipeline, for example, combines this approach
with calibrated sky models and empirical corrections \citep{Guy2023}.

Correlations among spectra further constrain the residual errors.
\citet{Kurtz2000} introduced eigenvector sky models derived by singular value
decomposition, and \citet{Wild2005} used principal component analysis (PCA)
to remove correlated OH residuals from SDSS spectra. These methods represent
structured residuals with a few basis spectra, whose fitted amplitudes refine
the initial sky estimate without a model for each residual feature.
\citet{Sharp2010} demonstrated substantial gains for deep fiber
observations, although the residuals decreased more slowly than Poisson
statistics would predict. Continuum subtraction also depends on how the sky
is sampled in space and on the observing strategy \citep{Yang2013}, and a
clean removal of the bright lines alone says little about the accuracy of the
underlying continuum. LAMOST combines a B-spline sky estimate, a sky-line
calibration and a PCA correction \citep{Bai2017}. Spatial gradients need
separate attention: \citet{Han2023} modelled the moonlight-induced variations
of intensity and colour across LAMOST fields with weighted gradient surfaces.

Instrument stability limits how accurately these estimates transfer between
fibers or exposures. Using MaNGA and VLT/FLAMES data, \citet{Bundy2022}
separated persistent biases from stochastic response variations caused by
changing fiber stress and the resulting focal ratio degradation. They found
that, with the errors under control, subtraction at about the
$0.1$--$0.2\%$ level is achievable.

Integral-field unit (IFU) spectroscopy extends these requirements to
spatially resolved emission from sources that may fill the science field,
which is why the MaNGA pipeline \citep{Law2016} and the SAMI reductions
\citep{Sharp2010,Sharp2015} use dedicated sky IFUs. MaNGA fits B-spline sky
models with adjustments for the background normalization and for
line-profile variations along the spectrograph slit, and reaches nearly
Poisson-limited subtraction shortward of about $8500\,\mathrm{\AA}$ in the
data examined \citep{Law2016}. SAMI combines 26 throughput-calibrated sky
fibers into a master spectrum for each exposure \citep{Sharp2015}, and
deliberately does not apply a PCA residual correction to science spectra,
because coherent galaxy features can be mistaken for sky structure. Cube
reconstruction introduces spatial covariance, and the uncertainty of a shared
sky estimate further limits what combining spectra can gain
\citep{Sharp2015,Law2016}. Both effects matter most for diffuse emission and
the outskirts of galaxies.

Other IFU pipelines take complementary approaches. The MUSE pipeline
integrates sky subtraction with instrument calibration and cube
reconstruction \citep{Weilbacher2020}, while the Zurich Atmosphere Purge
isolates sky residuals for a PCA correction by filtering and spectral
segmentation \citep{Soto2016}. Set against SAMI, these pipelines show that
empirical subtraction succeeds only if it keeps source information out of the
fitted sky basis. Source masks, training spectra and the number of retained
components must therefore preserve astronomical features, particularly when
those features recur across many spatial elements.

Physical constraints can also reduce the freedom a sky model needs. In
near-infrared spectroscopy, \citet{Davies2007} modeled variations among
groups of OH lines and small wavelength shifts, and \texttt{skycorr} extends
physically motivated line-group scaling across optical and near-infrared
spectra, fitting the lines separately from the continuum \citep{Noll2014}.
More recently, PALACE combines spectroscopic data with extensive Paranal
observations to model airglow lines, unresolved continuum components and
their climatological variability \citep{Noll2025}. Airglow intensities vary
on several timescales, so typical behavior and a particular exposure can
differ considerably. Such predictions can inform a fit to one exposure, but
they still have to be adapted to the instantaneous sky and instrumental
response, and other background components, such as scattered moonlight, need
a treatment of their own.

Recent data-driven methods complement these physical models.
\citet{Kolganov2025} used non-negative matrix factorization to model
two-dimensional optical sky spectra, including short-slit observations that
the source fills, without separate offset-sky exposures. In the
near-infrared, \citet{Dauphin2024} combined PCA reconstruction and
Gaussian-process analysis to characterize OH variability and moonlight, with
implications for observing cadence and sky modelling. These developments
point toward combining instrumental forward models, physical sky components
and learned spectral or temporal correlations. Success, in the end, means
recovering faint source signals with reliable uncertainties, as well as
reducing the sky residuals.
